\chapter{Reflection}\label{ch:reflection}

As I have described in \autoref{chapter:planning} it is really important for a game engine to have a reflection system that allows us to save and load map files and also edit the fields in a class from a nice graphical interface. While this is a really useful ability, \Cpp doesn't have it because it is not needed in every case and when it is not needed it would use up resources unnecessarily.

Because of this there are many different solutions to creating a reflection system on top of \Cpp. The easiest way is to create specialized template functions that give back the necessary information. For smaller projects or cases where new classes get added less frequently making there by hand is a good solution.

\begin{listing}[!ht]
    \begin{minted}{cpp}
    class Base {
    public:
        virtual const MetaData& getMetaData() const = 0;
    };

    class Derived : public Base {
    public:
        const MetaData& getMetaData() const override {
            static MetaData data = {
                .name = "Derived",
                .fields = {
                    { "field1", ... },
                    { "field2", ... },
                    { "field3", ... },
                }
            };
            return data;
        }
    };
    \end{minted}
    \caption{Example for a manual reflection function}
    \label{listing:maual-reflection}
\end{listing}

In the example (\autoref{listing:maual-reflection}), we have a base class that has a virtual function that all derived classes must overwrite and return the class metadata.
Using this solution we can store and get back the information we need in any format thus making it really flexible.
However, this is a lot of manual work and puts the responsibility on the coder to write the function correctly.
This takes time away from actual development.

Another solution is to use a tool that can generate these functions or generate the data in a portable format that can be loaded during runtime.


\section{Other engines}

In this section I will be presenting how other engines solve this problem by presenting two of them. I will not be discussing Unity here as it is not open source and uses \Csh 's built in reflection to some extent.

\subsection{Source}
Source uses the simplest approach where the developer is responsible for creating a function that returns all the metadata about a class.
To make this easier they have made use of the pre-processor macros in order to somewhat generalize the necessary lines of code. An example of this can be see in \autoref{listing:source-reflection} where a simple light's possible data fields are defined.

As source uses an external editor that doesn't depend on the game source code they also create a description file from these that tells the Hammer editor what entities exist in the game with what possible properties. This capability is useful if the engine supports outside editors for creating the levels.

\begin{listing}
    \begin{minted}{cpp}
BEGIN_DATADESC( CLight )

	DEFINE_FIELD( m_iCurrentFade, FIELD_CHARACTER),
	DEFINE_FIELD( m_iTargetFade, FIELD_CHARACTER),

	DEFINE_KEYFIELD( m_iStyle, FIELD_INTEGER, "style" ),
	DEFINE_KEYFIELD( m_iDefaultStyle, FIELD_INTEGER, "defaultstyle" ),
	DEFINE_KEYFIELD( m_iszPattern, FIELD_STRING, "pattern" ),

	// Fuctions
	DEFINE_FUNCTION( FadeThink ),

	// Inputs
	DEFINE_INPUTFUNC( FIELD_STRING, "SetPattern", InputSetPattern ),
	DEFINE_INPUTFUNC( FIELD_STRING, "FadeToPattern", InputFadeToPattern ),
	DEFINE_INPUTFUNC( FIELD_VOID,	"Toggle", InputToggle ),
	DEFINE_INPUTFUNC( FIELD_VOID,	"TurnOn", InputTurnOn ),
	DEFINE_INPUTFUNC( FIELD_VOID,	"TurnOff", InputTurnOff ),

END_DATADESC()
    \end{minted}
    \caption{Example for a manual reflection function}
    \label{listing:source-reflection}
\end{listing}

\subsection{Unreal Engine}

UE uses a separate build tool to generate reflection information by reading in an parsing the \Cpp source files of the project. This tool is strongly integrated into the build process and generates multiple extra source files for each class. It generates a new \Cpp and header file for each class that contains the reflection information. To know witch fields and in what format to represent the tool uses multiple special pre-processor macros that act like decorators or attributes in other languages. For example every class that wants to have reflection generated for it needs to start with \verb|UCLASS()| to make the tool aware of it. Similarly for each function and property the \verb|UFUNCTION()|, and \verb|UPROPERTY()| macros have to be placed before them to be recognised. An example of this can be seen in \autoref{listing:ue5-pawn-class} and some of the resulting generated code in \autoref{listing:ue5-generated-reflection}.


\begin{listing}
    \begin{minted}{cpp}
// Copyright Epic Games, Inc. All Rights Reserved.
#pragma once
#include "LyraPawn.generated.h"
UCLASS()
class LYRAGAME_API ALyraPawn : public AModularPawn, public ILyraTeamAgentInterface
{
	GENERATED_BODY()
private:
	UPROPERTY(ReplicatedUsing = OnRep_MyTeamID)
	FGenericTeamId MyTeamID;

private:
	UFUNCTION()
	void OnRep_MyTeamID(FGenericTeamId OldTeamID);
};
    \end{minted}
    \caption{Sections from a UE5 pawn class (Lyra demo project)}
    \label{listing:ue5-pawn-class}
\end{listing}

\begin{listing}
    \begin{minted}{cpp}
const FStructPropertyParams [..]_ALyraPawn_Statics::NewProp_MyTeamID = {
"MyTeamID",
"OnRep_MyTeamID",
(EPropertyFlags)0x0040000100000020,
UECodeGen_Private::EPropertyGenFlags::Struct,
RF_Public|RF_Transient|RF_MarkAsNative,
1,
nullptr,
nullptr,
STRUCT_OFFSET(ALyraPawn, MyTeamID),
Z_Construct_UScriptStruct_FGenericTeamId,
[...]
};
    \end{minted}
    \caption{Section from the generated reflection code for the same class as in \autoref{listing:ue5-pawn-class}}
    \label{listing:ue5-generated-reflection}
\end{listing}


\section{Shadow header tool}\label{sec:shadow-header-tool}

For Shadow Engine I have decided to try and create a similar tool to what UE has.
This would allow the users of the project a more seamless experience and free up development time in the future.
I have explored multiple possible ways to create such a tool, as for this I would need a partial AST (Abstract syntax tree).
This is notoriously hard to generate for \Cpp.

I have found multiple online tutorials on how to create compiler plugins for clang.
Clang is a compiler frontend for the famous LLVM compiler framework, and is one of the 3 most used C/\Cpp compilers.
While these tutorials showed great results and allowed full integration into the compilation process, they almost all required compiling clang from scratch.

I have tried following some of these tutorials, but I have run into many issues that would have need way more time to solve than I had available.
Most of my problems boiled down to one thing, there is no distribution of clang for Windows that supports loading and developing plugins.
Because of this I would need to compile clang myself and this is not an easy task and can differ based on the system being used.

After realising that getting clang plugins working is a way bigger task than I should work on, I have decided to write a simple parser in \Csh.
I have chosen \Csh because I'm familiar with it and has better string manipulation methods than cpp.

...

Cpp attributes

\subsection{Goal}
Get a list of all the classes that are in the code.
To be expandable in the future the tool parses all the classes no matter if they have an attribute that marks it for reflection.

\subsection{Ast}
After the initial regex based implementation I got inspiration and initial help creating a more advanced code processor that is much closerin design to a classical compiler.
These software usually have the following steps:
- Lexer
- parser
- etc..
The lexer is responsible for dividing the source code into lexicographical tokens.
These are usually things like keywords, identifiers, and the different kinds of brackets.
This is done by iterating over the input text character by character and if the character is a known char we put it into a list of tokens.
<Example code>
In my implementation I parse the following tokens:
<List>
These are not all the possible tokens, but they are enough for what this tool needs to understand.

\subsection{Parser}
Parsing is the step where we iterate over all the tokens that makeup a file and form logical units out of them.
These are going to be the Nodes that make up the final AST. For each file the root of the corresponding AST is going to be a CompilationUnitNode.

For example the following series of tokens represent a basic class.

<Example of class tokens>

Would turn into this:

<Pseudo code of class node>

During parsing we find these patterns and form node from them.

To get this we can do a similar forward iteration as we did while lexing but this time we iterate on the token level.
Most of these patterns are well documented in the cop specifications and are more easily available online in cppreference.com

In the tool I have created a partial parser as I do not care about most of the cop code.
As a result my parser ignores all the function bodies and any declarations that it doesn't recognise as a class head or a member field.

\subsection{Output}
Both of the parsers (regex based and lexing) output the data they gather in a common format.
This makes it easy to add new code processors in the future, thus leaving the possibility of a clang integration open.

The output writer is a simple function that finds all classes with the correct attribute (SH::Reflect) and,
using simple string interpolation writes the necessary functions into a common file.
